\section{Order is a different structure}
Equivalence relations describe sameness. A different kind of relation describes \emph{comparison}: one object is smaller than, below, or contained in another. Write $x\preceq y$ for such a relation, read ``$x$ is below $y$.'' It is a \emph{partial order} when it satisfies three requirements for all allowed $x,y,z$:
\begin{itemize}
\item \emph{Reflexivity:} $x\preceq x$.
\item \emph{Antisymmetry:} if $x\preceq y$ and $y\preceq x$, then $x=y$.
\item \emph{Transitivity:} if $x\preceq y$ and $y\preceq z$, then $x\preceq z$.
\end{itemize}
Reflexivity and transitivity are the same requirements as for an equivalence relation, but symmetry has been replaced by antisymmetry. Symmetry says that a relation always goes both ways; antisymmetry says that it goes both ways only between an object and itself. The usual order $\le$ on the real numbers is a partial order.

Set inclusion $\subseteq$ is also a partial order on any collection of sets. Unlike $\le$ on numbers, it need not compare every pair: neither $\{1\}$ nor $\{2\}$ contains the other. That is why the word \emph{partial} appears in the name.

Because some pairs may be incomparable, there are two different notions of a ``largest'' element of a collection, and they must not be confused.
\begin{itemize}
\item An element $x$ is \emph{maximal} if nothing in the collection lies strictly above it: there is no $y$ in the collection with $x\preceq y$ and $y\ne x$.
\item An element $x$ is a \emph{maximum} if everything in the collection lies below it: $y\preceq x$ for every $y$ in the collection.
\end{itemize}
A maximum is always maximal, but not conversely. In the collection $\{\{1\},\{2\}\}$, ordered by inclusion, both members are maximal, because neither lies inside the other. Neither member is a maximum, because neither contains the other. A collection can therefore have several maximal elements and no maximum.

A related but different distinction appears in Chapter~8. There, a partial assignment of tasks can be \emph{maximal under inclusion}---no single new pair can be added to it---without having the largest possible \emph{number} of pairs. In that chapter the word ``maximum'' will mean ``largest number of pairs,'' not ``contains every other assignment''; two assignments of the largest size need not contain one another. In both settings the lesson is the same: being impossible to improve by one obvious step is weaker than being best overall.

\begin{principle}{Habit 6: ask what a change of description preserves}
Before quotienting, normalizing, or identifying objects, specify the relation and the preserved property. Before applying a formula to equivalence classes, check independence of representatives. A short well-definedness check can prevent an entire argument from being about the wrong object.
\end{principle}
